\def\curriculumtitle{Chapter 1: What a Number Means}
%% shared preamble for every chapter of the teaching PDF curriculum
%% each chapter imports this via %% shared preamble for every chapter of the teaching PDF curriculum
%% each chapter imports this via %% shared preamble for every chapter of the teaching PDF curriculum
%% each chapter imports this via \input{../../chapter_preamble.tex}
%% style adapted from neuroloc_guide.tex

\documentclass[11pt,a4paper]{article}
\usepackage[top=2.5cm,bottom=2.5cm,left=2.5cm,right=2.5cm]{geometry}
\usepackage{amsmath,amssymb,amsthm,mathtools}
\usepackage{bm}
\usepackage{enumitem}
\usepackage{titlesec}
\usepackage{fancyhdr}
\usepackage{parskip}
\usepackage[T1]{fontenc}
\usepackage{lmodern}
\usepackage{microtype}
\usepackage{hyperref}
\hypersetup{hidelinks}
\usepackage{tikz}
\usepackage{xcolor}

\pagestyle{fancy}
\fancyhf{}
\fancyhead[L]{\small\textit{\curriculumtitle}}
\fancyhead[R]{\small\thepage}
\renewcommand{\headrulewidth}{0.4pt}

\titleformat{\section}{\Large\bfseries}{}{0em}{}
\titleformat{\subsection}{\large\bfseries}{}{0em}{}
\titleformat{\subsubsection}{\normalsize\bfseries}{}{0em}{}

\setlist[itemize]{topsep=4pt,itemsep=2pt,parsep=0pt}
\setlist[enumerate]{topsep=4pt,itemsep=2pt,parsep=0pt}

\newtheoremstyle{ch1-defn}{6pt}{6pt}{}{}{\bfseries}{.}{0.5em}{}
\newtheoremstyle{ch1-example}{6pt}{6pt}{}{}{\bfseries}{.}{0.5em}{}

\theoremstyle{ch1-defn}
\newtheorem{definition}{Definition}[section]
\newtheorem{lemma}[definition]{Lemma}
\newtheorem{theorem}[definition]{Theorem}

\theoremstyle{ch1-example}
\newtheorem{example}[definition]{Example}
\newtheorem{remark}[definition]{Remark}

%% callout box for "common mistakes" in the dzipobel pattern
\newenvironment{mistake}{%
  \begin{quote}\small\textbf{Common mistake.}\ }{\end{quote}}

%% callout box for "verification trick" (minimal-pair mnemonic)
\newenvironment{trick}{%
  \begin{quote}\small\textbf{Verification trick.}\ }{\end{quote}}

%% default curriculum title, override per-chapter if desired
\providecommand{\curriculumtitle}{Todorov Teaching Curriculum}

%% style adapted from neuroloc_guide.tex

\documentclass[11pt,a4paper]{article}
\usepackage[top=2.5cm,bottom=2.5cm,left=2.5cm,right=2.5cm]{geometry}
\usepackage{amsmath,amssymb,amsthm,mathtools}
\usepackage{bm}
\usepackage{enumitem}
\usepackage{titlesec}
\usepackage{fancyhdr}
\usepackage{parskip}
\usepackage[T1]{fontenc}
\usepackage{lmodern}
\usepackage{microtype}
\usepackage{hyperref}
\hypersetup{hidelinks}
\usepackage{tikz}
\usepackage{xcolor}

\pagestyle{fancy}
\fancyhf{}
\fancyhead[L]{\small\textit{\curriculumtitle}}
\fancyhead[R]{\small\thepage}
\renewcommand{\headrulewidth}{0.4pt}

\titleformat{\section}{\Large\bfseries}{}{0em}{}
\titleformat{\subsection}{\large\bfseries}{}{0em}{}
\titleformat{\subsubsection}{\normalsize\bfseries}{}{0em}{}

\setlist[itemize]{topsep=4pt,itemsep=2pt,parsep=0pt}
\setlist[enumerate]{topsep=4pt,itemsep=2pt,parsep=0pt}

\newtheoremstyle{ch1-defn}{6pt}{6pt}{}{}{\bfseries}{.}{0.5em}{}
\newtheoremstyle{ch1-example}{6pt}{6pt}{}{}{\bfseries}{.}{0.5em}{}

\theoremstyle{ch1-defn}
\newtheorem{definition}{Definition}[section]
\newtheorem{lemma}[definition]{Lemma}
\newtheorem{theorem}[definition]{Theorem}

\theoremstyle{ch1-example}
\newtheorem{example}[definition]{Example}
\newtheorem{remark}[definition]{Remark}

%% callout box for "common mistakes" in the dzipobel pattern
\newenvironment{mistake}{%
  \begin{quote}\small\textbf{Common mistake.}\ }{\end{quote}}

%% callout box for "verification trick" (minimal-pair mnemonic)
\newenvironment{trick}{%
  \begin{quote}\small\textbf{Verification trick.}\ }{\end{quote}}

%% default curriculum title, override per-chapter if desired
\providecommand{\curriculumtitle}{Todorov Teaching Curriculum}

%% style adapted from neuroloc_guide.tex

\documentclass[11pt,a4paper]{article}
\usepackage[top=2.5cm,bottom=2.5cm,left=2.5cm,right=2.5cm]{geometry}
\usepackage{amsmath,amssymb,amsthm,mathtools}
\usepackage{bm}
\usepackage{enumitem}
\usepackage{titlesec}
\usepackage{fancyhdr}
\usepackage{parskip}
\usepackage[T1]{fontenc}
\usepackage{lmodern}
\usepackage{microtype}
\usepackage{hyperref}
\hypersetup{hidelinks}
\usepackage{tikz}
\usepackage{xcolor}

\pagestyle{fancy}
\fancyhf{}
\fancyhead[L]{\small\textit{\curriculumtitle}}
\fancyhead[R]{\small\thepage}
\renewcommand{\headrulewidth}{0.4pt}

\titleformat{\section}{\Large\bfseries}{}{0em}{}
\titleformat{\subsection}{\large\bfseries}{}{0em}{}
\titleformat{\subsubsection}{\normalsize\bfseries}{}{0em}{}

\setlist[itemize]{topsep=4pt,itemsep=2pt,parsep=0pt}
\setlist[enumerate]{topsep=4pt,itemsep=2pt,parsep=0pt}

\newtheoremstyle{ch1-defn}{6pt}{6pt}{}{}{\bfseries}{.}{0.5em}{}
\newtheoremstyle{ch1-example}{6pt}{6pt}{}{}{\bfseries}{.}{0.5em}{}

\theoremstyle{ch1-defn}
\newtheorem{definition}{Definition}[section]
\newtheorem{lemma}[definition]{Lemma}
\newtheorem{theorem}[definition]{Theorem}

\theoremstyle{ch1-example}
\newtheorem{example}[definition]{Example}
\newtheorem{remark}[definition]{Remark}

%% callout box for "common mistakes" in the dzipobel pattern
\newenvironment{mistake}{%
  \begin{quote}\small\textbf{Common mistake.}\ }{\end{quote}}

%% callout box for "verification trick" (minimal-pair mnemonic)
\newenvironment{trick}{%
  \begin{quote}\small\textbf{Verification trick.}\ }{\end{quote}}

%% default curriculum title, override per-chapter if desired
\providecommand{\curriculumtitle}{Todorov Teaching Curriculum}


\begin{document}

%% ============================================================
%%  COVER
%% ============================================================

\thispagestyle{empty}
\vspace*{5cm}
\begin{center}
{\Huge\bfseries Chapter 1}\\[0.4cm]
{\Huge\bfseries What a Number Means}\\[1.5cm]
{\Large A Teaching Curriculum from First Principles}\\[0.4cm]
{\normalsize Prerequisites: none}\\[2cm]
{\large Deyan Todorov}\\[0.2cm]
{\large Eptesicus Labs}\\[1.5cm]
{\normalsize 2026}
\end{center}
\newpage

%% ============================================================
%%  OPENING
%% ============================================================

\section*{Opening}

This chapter answers one question: \textit{what is a number, and why do we use symbols to stand in for them?}

You already use numbers every day. You count apples. You measure rooms. You read phone numbers. But when mathematicians write $\mathbb{R}$ or $f : A \to B$ or $\forall x \in \mathbb{N}$, they mean something precise. The precision matters. Every chapter after this one uses numbers and symbols in that precise sense, and if you try to learn derivatives (Chapter~2) or probability (Chapter~7) or attention (Chapter~26) while treating symbols as something you pattern-match on instead of something you understand, you will hit a wall.

We build up from absolute zero prerequisites. Seven sections, one worked example, one quiz:

\begin{enumerate}
  \item What ``number'' refers to.
  \item Counting with the integers.
  \item Measuring with the rationals and reals.
  \item Symbols and variables.
  \item Functions.
  \item Domain, codomain, range.
  \item Injective, surjective, bijective.
\end{enumerate}

\noindent\textbf{Convention.} Throughout this chapter we use $\mathbb{N} = \{0, 1, 2, 3, \ldots\}$ (the natural numbers include zero) and $\mathbb{Z}^{+} = \{1, 2, 3, \ldots\}$ (the strictly positive integers). Some older texts write $\mathbb{N} = \{1, 2, 3, \ldots\}$ instead. Both conventions coexist in the literature. We pick the ISO/computing convention because induction starts at~$0$ and array indices start at~$0$, and we want the symbol $\mathbb{N}$ to match that.

%% ============================================================
%%  SECTION 1
%% ============================================================

\section{What ``number'' refers to}

The word ``number'' in everyday English points at four different things. Only two of them are numbers in the mathematical sense of this chapter.

\begin{enumerate}
  \item \textbf{Counts.} ``Three apples.'' The number three here answers \textit{how many}. It can be added (three apples plus two apples is five apples), subtracted, compared (three is less than five).
  \item \textbf{Measurements.} ``3.5 meters.'' The number answers \textit{how much}. It participates in arithmetic the same way counts do, just with finer resolution.
  \item \textbf{Labels.} ``Phone number 555-1234.'' The digits identify a specific phone, but the expression is not arithmetic: phone 555-1234 plus phone 555-5678 is nonsense. A label happens to use digits; that does not make it a mathematical number.
  \item \textbf{Positions.} ``Third in line.'' The digit identifies a location in an order. ``Third plus fifth'' is meaningful only if you interpret it as a count (the eighth position) or as positions in a different sequence.
\end{enumerate}

Mathematical numbers are counts and measurements. The defining property: they participate in arithmetic and comparison. You can add them, subtract them, multiply them, compare them. Labels and positions use digits but do not participate the same way.

\begin{mistake}
``Phone number 555-1234 is a number.''\ Not in the mathematical sense. It is a \emph{label} that uses digits. A mathematical number is an object in $\mathbb{N}$, $\mathbb{Z}$, $\mathbb{Q}$, or $\mathbb{R}$ that can be added, subtracted, multiplied, and compared. Digits can appear in non-numerical roles: phone numbers, house numbers, jersey numbers, passport numbers. None of these participate in arithmetic.
\end{mistake}

The distinction looks trivial. It is load-bearing. Later in the curriculum, when we talk about ``indexing into a tensor'' (Chapter~4) or ``the \(k\)-th element of a sequence'' (Chapter~25), we will be using numbers as \emph{positions}, not as counts. And when we measure the output of a probability distribution (Chapter~7), we will be using numbers as \emph{measurements}. Knowing which one you mean prevents category errors that are painful to untangle later.

\subsection*{Worked example}

Which of the following four expressions are mathematical numbers, and which are not?

\begin{enumerate}
  \item \textbf{Three apples.} Mathematical number. The number~$3$ is a count; it participates in arithmetic.
  \item \textbf{Phone number 555-1234.} Not a mathematical number. It is a label.
  \item \textbf{The third room on the left.} The number~$3$ here is a \emph{position} (third). Treated as a count of positions from the start, it is mathematical; treated as a standalone label ``room~3'', it is not. Context decides.
  \item \textbf{3.5 meters.} Mathematical number. It is a measurement.
\end{enumerate}

The lesson: mathematical numbers are the ones that survive being added, compared, or used in a calculation. When in doubt, try to add them to something. If the addition is nonsense, you are not holding a mathematical number.

%% ============================================================
%%  SECTION 2
%% ============================================================

\section{Counting with the integers}

Start with the simplest counts: zero, one, two, three, and so on. The set of all such counts is the \textbf{natural numbers}:
\[
  \mathbb{N} = \{0, 1, 2, 3, 4, \ldots\}.
\]
A \textbf{set}, in the informal sense we use here, is just an unordered collection of distinct objects. The braces denote membership. ``$5 \in \mathbb{N}$'' reads ``$5$ is an element of $\mathbb{N}$''; it is true. ``$-1 \in \mathbb{N}$'' is false, because $-1$ is not in the collection.

\subsection*{What the naturals can and cannot do}

Pick any two naturals and add them: $3 + 4 = 7$. The result is a natural. Do the same with multiplication: $3 \times 4 = 12$. Also a natural. We say $\mathbb{N}$ is \textbf{closed under addition} and \textbf{closed under multiplication}: the result of the operation stays inside the set you started from.

Now try subtraction. $5 - 3 = 2$, a natural. Fine. But $3 - 5 = \,?\,$ There is no natural number that, when added to $5$, gives $3$. Subtraction of naturals sometimes has no answer inside $\mathbb{N}$. The set is \textbf{not closed under subtraction}.

This is not a small problem. If we want every subtraction to produce an answer, we have to extend the number system. The extension is the \textbf{integers}:
\[
  \mathbb{Z} = \{\ldots, -3, -2, -1, 0, 1, 2, 3, \ldots\}.
\]
The new entries are the negatives. Once we have them, $3 - 5 = -2$ is defined. The integers are closed under addition, subtraction, and multiplication.

\begin{example}[the motivating gap from $\mathbb{N}$ to $\mathbb{Z}$]
In $\mathbb{N}$: the expression $7 - 10$ has no answer. No natural number, when added to $10$, gives $7$. This gap is the reason we build $\mathbb{Z}$. Once in $\mathbb{Z}$, $7 - 10 = -3$ and the expression is well-defined.
\end{example}

\subsection*{But the integers have their own gap}

Division does not work inside $\mathbb{Z}$. $10 \div 5 = 2$, that is fine. But $7 \div 10 = \,?\,$ There is no integer that, when multiplied by $10$, gives $7$. The integers are not closed under division.

This forces another extension, to the \textbf{rational numbers}:
\[
  \mathbb{Q} = \left\{\frac{p}{q} \,:\, p \in \mathbb{Z},\, q \in \mathbb{Z},\, q \neq 0\right\}.
\]
A rational is a ratio of integers with nonzero denominator. $7/10$, $-3/4$, $22/7$, and $5$ (which is the ratio $5/1$) are all rationals. The rationals are closed under addition, subtraction, multiplication, and division by anything nonzero.

\begin{example}[the motivating gap from $\mathbb{Z}$ to $\mathbb{Q}$]
In $\mathbb{Z}$: the expression $7 \div 10$ has no answer. No integer, when multiplied by $10$, gives $7$. Once in $\mathbb{Q}$, $7 \div 10 = 7/10$ and the expression is well-defined.
\end{example}

\subsection*{The pattern}

Each extension is driven by a specific operation that had no answer in the previous system. We extend only as much as needed. That is the whole logic of the ladder:
\[
  \mathbb{N} \;\subset\; \mathbb{Z} \;\subset\; \mathbb{Q} \;\subset\; \mathbb{R}.
\]
We still have one more rung to go. Division closed $\mathbb{Q}$ under the four basic operations. What remains to fail in $\mathbb{Q}$ is something more subtle, and it is what forces us to the reals. That is the subject of the next section.

\begin{remark}
Peano's axioms give a rigorous construction of $\mathbb{N}$ from set-theoretic primitives. We do not use them in this chapter. The ``closure under operations'' framing is enough for our purposes; Peano's axioms belong to a later chapter on logic and foundations.
\end{remark}

%% ============================================================
%%  SECTION 3
%% ============================================================

\section{Measuring with the rationals and reals}

The rationals have a useful property: between any two distinct rationals, there is another rational. Pick $1/2$ and $1/3$. Their average $\tfrac{1}{2}(1/2 + 1/3) = 5/12$ sits strictly between them. You can repeat the averaging as many times as you want, and every result stays rational. The rationals are \textbf{dense}: no matter how close you squeeze two of them, there is always another one in the gap.

Given how tightly packed they are, you might expect the rationals to cover the whole number line. They do not. They have holes.

\subsection*{The hole}

Consider the set of positive rationals whose square is less than~$2$:
\[
  C \;=\; \{x \in \mathbb{Q} \,:\, x > 0 \text{ and } x^2 < 2\}.
\]
Every element of $C$ is bounded above by $3/2$, because $(3/2)^2 = 9/4 > 2$. So $C$ has rational upper bounds. Does $C$ have a \emph{least} rational upper bound?

If a rational $b$ were the least upper bound of $C$, then $b^2$ would have to equal~$2$ exactly. Any $b$ with $b^2 < 2$ would not be an upper bound (you can find a slightly larger rational still in $C$), and any $b$ with $b^2 > 2$ would not be \emph{least} (you can find a slightly smaller rational still above~$C$). So the least upper bound, if it exists as a rational, must satisfy $b^2 = 2$.

But no rational satisfies $b^2 = 2$. We prove this next.

\subsection*{$\sqrt{2}$ is not rational}

This is the canonical proof by contradiction.

\begin{lemma}[used twice in the proof below]
If $k$ is an integer and $k^2$ is even, then $k$ is even.
\end{lemma}

\begin{proof}
Suppose $k$ is odd. Then $k = 2j + 1$ for some integer $j$, and
\[
  k^2 \;=\; (2j+1)^2 \;=\; 4j^2 + 4j + 1 \;=\; 2(2j^2 + 2j) + 1,
\]
which is odd. Contradiction. So if $k^2$ is even, $k$ cannot be odd. Therefore $k$ is even.
\end{proof}

\begin{theorem}[no positive rational squares to 2]
There is no positive rational number whose square is~$2$.
\end{theorem}

\begin{proof}
Suppose, for contradiction, that there were. Write it in lowest form as $p/q$, where $p$ and $q$ are positive integers with $\gcd(p, q) = 1$ (no common factor other than~$1$).

Then $(p/q)^2 = 2$, so
\[
  p^2 \;=\; 2q^2.
\]
The right side is twice an integer, so $p^2$ is even. By the lemma, $p$ is even. Write $p = 2k$ for some positive integer $k$. Substitute:
\[
  (2k)^2 \;=\; 2q^2 \quad\Rightarrow\quad 4k^2 \;=\; 2q^2 \quad\Rightarrow\quad q^2 \;=\; 2k^2.
\]
Now $q^2$ is even, so by the lemma again, $q$ is even.

But we assumed $\gcd(p, q) = 1$. Both $p$ and $q$ being even contradicts that. Therefore no rational $p/q$ squares to~$2$.
\end{proof}

\subsection*{The reals fill the hole}

The rationals cannot answer every question about least upper bounds. To fix this, we extend one more time, to the \textbf{real numbers} $\mathbb{R}$. The reals are defined to include every rational AND every ``hole-filler'' like $\sqrt{2}$, $\pi$, and $e$. The defining property of $\mathbb{R}$ (Spivak's axiom P13) is:

\begin{quote}
If $A \subseteq \mathbb{R}$ is a nonempty set with an upper bound, then $A$ has a least upper bound in $\mathbb{R}$.
\end{quote}

This is called the \textbf{completeness} of the reals. It is the one property that distinguishes $\mathbb{R}$ from $\mathbb{Q}$. Everything else about arithmetic (commutativity, distributivity, ordering) the rationals already have. Completeness is the new thing.

A number in $\mathbb{R}$ that is not in $\mathbb{Q}$ is called \textbf{irrational}. $\sqrt{2}$ is irrational. So are $\pi$ and $e$. Irrationals are real numbers; they live on the number line; they participate in arithmetic. The word ``irrational'' only means ``not rational.'' It does not mean ``not a number.''

\begin{mistake}
``$\sqrt{2}$ is a strange number that does not exist.''\ $\sqrt{2}$ is a real number with a definite location on the number line, approximately $1.4142135\ldots$. It is not \emph{rational} — it cannot be written as a ratio of two integers — but it is a number, and it participates in arithmetic like any other real. ``Irrational'' means \textit{not in $\mathbb{Q}$}; it does not mean \textit{not real} or \textit{not a number}.
\end{mistake}

\subsection*{A brief note on infinity}

Both $\mathbb{N}$ and $\mathbb{R}$ are infinite sets. But they are infinite in \emph{different ways}, and this distinction matters later in the curriculum.

You can list the naturals one at a time: $0, 1, 2, 3, \ldots$ and every natural eventually appears. Whatever natural you name, it shows up at a specific position in the list.

You cannot list the reals that way. Between any two reals there are uncountably many more reals, and no enumeration — no matter how clever — captures them all. A later chapter (on sets and cardinalities) proves this formally via Cantor's diagonal argument. For now, remember the shape:
\begin{itemize}
  \item $\mathbb{N}$ is infinite, but \emph{listable}.
  \item $\mathbb{R}$ is infinite, but \emph{not listable} in that way.
\end{itemize}
``Infinite'' is not one size. There are at least two sizes, and $\mathbb{R}$ is the bigger one.

%% ============================================================
%%  SECTION 4
%% ============================================================

\section{Symbols and variables}

A \textbf{variable} is a symbol (almost always a letter) that stands in for a number. When we write $x + 1 = 3$, we are not stating that the letter $x$ is equal to something. We are stating a relationship, and the letter $x$ is a placeholder for whatever value makes the relationship true.

Three things to keep separate:

\begin{enumerate}
  \item \textbf{The symbol.} The letter $x$ as an abstract name, ranging over some set (often $\mathbb{R}$).
  \item \textbf{A specific value of the symbol.} When we say ``let $x = 3$'', we pin the variable to a specific number. The letter no longer ranges; it has a value.
  \item \textbf{A proposition involving the symbol.} A statement like ``$x > 0$'' is a \emph{claim} that becomes true or false depending on the value of $x$.
\end{enumerate}

\subsection*{Equality is not assignment}

In a programming language, you may be used to writing \texttt{x = 3} to mean ``make the variable \texttt{x} hold the value $3$.'' That is \textbf{assignment}. In mathematics, the symbol $=$ means something different. It means \textbf{equality}: the two sides are claimed to refer to the same number.

\begin{mistake}
``$x + 1 = 3$ is an instruction to set $x + 1$ equal to $3$.''\ Not in mathematics. The symbol $=$ in math is a claim, not an instruction. The equation $x + 1 = 3$ is the claim ``whatever value $x$ has, $x + 1$ is the same number as~$3$.'' Your job as the reader is to figure out which values of $x$ make the claim true. In programming, the operator \(:=\) or \(\leftarrow\) is often used for assignment specifically to avoid this confusion; in math, $=$ is \emph{always} equality.
\end{mistake}

\subsection*{Equation vs.\ identity}

Two sentences that look the same can do very different jobs.

\begin{trick}
Read each of these and decide whether the equals sign is a \emph{question} (``for which values of $x$?'') or a \emph{statement} (``for every $x$.''):
\begin{itemize}
  \item $x + 1 = 3$. A \textbf{conditional equation}. True for one specific value: $x = 2$. Your job is to solve for that value.
  \item $x + 1 = 3x$. Also a conditional equation. True for $x = 1/2$. Solve for it.
  \item $x + 1 = 1 + x$. An \textbf{identity}. True for \emph{every} value of $x$. There is nothing to solve. The two sides are the same object written two ways.
\end{itemize}
\end{trick}

The difference is not cosmetic. Later chapters depend on it: when we say ``$\sin^2 x + \cos^2 x = 1$'' we mean an identity (holds for all $x$), but when we say ``$\sin x = 0$'' we mean a conditional equation (holds for $x = 0, \pm\pi, \pm 2\pi, \ldots$). Misreading one as the other produces wrong algebra.

\subsection*{A preview of quantifiers}

When you want to say ``for every $x$, something is true,'' mathematicians write
\[
  \forall x \colon P(x),
\]
pronounced ``for all $x$, $P$ of $x$.'' The symbol $\forall$ is the \textbf{universal quantifier}.

When you want to say ``there exists some $x$ for which something is true,'' mathematicians write
\[
  \exists x \colon P(x),
\]
pronounced ``there exists $x$ such that $P$ of $x$.'' The symbol $\exists$ is the \textbf{existential quantifier}.

Two examples:
\begin{align*}
  \forall x \in \mathbb{R}, \;\; x^2 &\geq 0 & &\text{(every real has a non-negative square)} \\
  \exists x \in \mathbb{R}, \;\; x^2 &= 4   & &\text{(some real squares to 4 --- in fact two of them)}
\end{align*}

This is a preview, not the full treatment. A later chapter develops formal quantifier reasoning. For now, you need to recognize the symbols and know what they assert.

%% ============================================================
%%  SECTION 5
%% ============================================================

\section{Functions}

A \textbf{function} is a rule. It takes an input from one set and produces an output in another set, with one condition: each input produces \emph{exactly one} output.

\begin{definition}[function, after Rosen]
Let $A$ and $B$ be nonempty sets. A function $f$ from $A$ to $B$ is an assignment of exactly one element of $B$ to each element of $A$. We write $f(a) = b$ if $b$ is the unique element of $B$ that $f$ assigns to $a$. We denote the whole thing by
\[
  f \colon A \to B.
\]
\end{definition}

Rosen states the definition only for nonempty sets. That restriction is useful here because it keeps us away from edge cases before we need them. In full set theory, mathematicians do allow functions whose domain is empty. We do not need that case in this chapter.

\subsection*{Four ways to specify a function (after Stewart)}

The same function can be described in four equivalent styles. A working mathematician moves between them constantly.

\begin{enumerate}
  \item \textbf{Verbal.} ``The rule that doubles every real number and adds one.''
  \item \textbf{Algebraic (formula).} $f(x) = 2x + 1$.
  \item \textbf{Graphical.} The set of points $(x, 2x + 1)$ drawn on the $xy$-plane. A straight line of slope $2$ crossing the $y$-axis at $(0, 1)$.
  \item \textbf{Numerical (table).}
    \[
      \begin{array}{c|c}
        x & f(x) \\ \hline
        0 & 1 \\ 1 & 3 \\ 2 & 5 \\ 3 & 7
      \end{array}
    \]
\end{enumerate}

No single representation is ``the'' function. The function is the rule; all four encode it.

\subsection*{The vertical line test}

A useful geometric consequence of the definition: in any graph that is the graph of a function, every vertical line crosses the graph at most once. Why? Because a vertical line $x = c$ corresponds to a single input value $c$, and the definition allows only one output. Two outputs would mean two intersection points.

This is worth stating as a criterion: a curve in the plane is the graph of a function of $x$ if and only if every vertical line crosses it at most once.

\subsection*{A minimal pair: function vs.\ relation (after Dawkins)}

Compare these two equations side by side:

\begin{align*}
  y &= x^2 + 1 \\
  y^2 &= x + 1
\end{align*}

The first defines a function. For every input $x$, there is exactly one output $y$. At $x = 3$, $y = 3^2 + 1 = 10$. Always one answer.

The second does not. At $x = 3$, we have $y^2 = 4$, so $y = +2$ or $y = -2$. Two possible outputs for the same input. It fails the ``exactly one'' rule; it fails the vertical line test (a vertical line at $x = 3$ crosses the curve twice); it is not a function. It is a more general object called a \textbf{relation}.

\begin{mistake}
``$y^2 = x + 1$ is a function because $y$ appears on the left.'' Not a function. The test is whether each input produces exactly one output, not which letter appears where. At $x = 3$, the equation has two solutions $y = \pm 2$.

Reminder (Dawkins): \textit{``This is NOT a multiplication of $f$ by $x$! This is one of the more common mistakes people make when they first deal with functions.''} The expression $f(x)$ is a function call. The parentheses do \emph{not} mean multiplication in this context.
\end{mistake}

\subsection*{The transformation framing (after 3Blue1Brown)}

Once the formal definition is in place, a useful intuition: a function is something that \emph{moves} each input to its output. If $f : \mathbb{R} \to \mathbb{R}$ is the rule $f(x) = 2x + 1$, then $f$ moves $0$ to $1$, moves $1$ to $3$, moves $2$ to $5$, and so on. The value $f(x)$ is \textit{where $x$ goes}.

This kinematic image is convenient later: composing two functions is ``first move, then move again''; classifying functions is ``movements that do not collide'' (injective) and ``movements that cover the whole destination'' (surjective). We will use these pictures in Section~7.

%% ============================================================
%%  SECTION 6
%% ============================================================

\section{Domain, codomain, range}

A function $f : A \to B$ has three associated sets. They are different, and the chapter will harp on the difference because it matters for Section~7.

\begin{definition}[the three sets]
Let $f : A \to B$ be a function.
\begin{itemize}
  \item The \textbf{domain} of $f$ is $A$: the set of allowed inputs.
  \item The \textbf{codomain} of $f$ is $B$: the \emph{declared} target set. Every output is guaranteed to live in $B$, but not every element of $B$ is guaranteed to be hit by some input.
  \item The \textbf{range} of $f$ is the set of outputs that actually occur:
    \[
      \mathrm{range}(f) \;=\; \{f(a) \,:\, a \in A\}.
    \]
    The range is always a subset of the codomain: $\mathrm{range}(f) \subseteq B$.
\end{itemize}
\end{definition}

\subsection*{Why codomain and range are not the same}

Consider $f : \mathbb{R} \to \mathbb{R}$ defined by $f(x) = x^2$.

\begin{itemize}
  \item The domain is $\mathbb{R}$. Every real number is allowed as an input.
  \item The codomain is $\mathbb{R}$. That is how we declared the function.
  \item The range is $[0, \infty)$. The outputs are all non-negative reals. No input produces a negative output, because $x^2 \geq 0$ for every real $x$.
\end{itemize}

Notice: $\mathrm{range}(f) = [0, \infty) \;\subsetneq\; \mathbb{R} = B$. The range is a \emph{proper} subset of the codomain. This is normal; for most functions, the codomain is a declared ``generous'' target and the range is the actual, possibly smaller, image.

\begin{mistake}
``The range of $f(x) = x^2$ from $\mathbb{R}$ to $\mathbb{R}$ is $\mathbb{R}$, because $\mathbb{R}$ is the target set written to the right of the arrow.''

That is the canonical conflation. It is what many calculus textbooks (Stewart among them) implicitly teach by collapsing codomain and range into one word. In this curriculum we keep them separate.

\textbf{Right.} The range of $f(x) = x^2$ from $\mathbb{R}$ to $\mathbb{R}$ is $[0, \infty)$. This is a \emph{proper subset} of the codomain $\mathbb{R}$. Negative numbers are in the codomain (we declared $\mathbb{R}$) but not in the range (no real input produces a negative $x^2$). Codomain is declared; range is actual; they coincide only when the function is surjective (Section~7).
\end{mistake}

\subsection*{Why this distinction will matter in a moment}

Section~7 introduces three classifications of functions: injective, surjective, and bijective. All three depend on \emph{which} sets we declared as domain and codomain, not just on the formula. Changing the codomain while keeping the formula can flip a function from not-surjective to surjective without touching a single value. Same formula, different classification.

That is why we insist on three named sets, not two.

%% ============================================================
%%  SECTION 7
%% ============================================================

\section{Injective, surjective, bijective}

Three questions you can ask about any function $f : A \to B$:

\begin{enumerate}
  \item Do different inputs always produce different outputs? If yes, $f$ is \textbf{injective} (also called \textbf{one-to-one}).
  \item Does every element of the codomain $B$ get hit by some input? If yes, $f$ is \textbf{surjective} (also called \textbf{onto}).
  \item Is $f$ both injective and surjective? If yes, $f$ is \textbf{bijective} (also called a \textbf{one-to-one correspondence}).
\end{enumerate}

We use the Latin terms (injective, surjective, bijective) and the English phrases (one-to-one, onto, one-to-one correspondence) interchangeably. Both are standard. The Latin terms are unambiguous; the English phrases have a naming collision you must watch for.

\begin{mistake}
``One-to-one'' means injective. ``One-to-one correspondence'' means bijective. The short phrase and the long phrase mean \emph{different things}. Do not treat them as synonyms.

If the English ambiguity ever bites you, fall back on the Latin. Injective means injective; bijective means bijective; there is no room for confusion.
\end{mistake}

\subsection*{Formal definitions}

\begin{definition}[injective]
$f : A \to B$ is \textbf{injective} if, whenever $f(a_1) = f(a_2)$ for some $a_1, a_2 \in A$, we have $a_1 = a_2$. Equivalently (by contrapositive): whenever $a_1 \neq a_2$, we have $f(a_1) \neq f(a_2)$.
\end{definition}

The two forms are logically equivalent. The direct form is easier when you already know the outputs are equal and want to conclude something about the inputs. The contrapositive is easier when you are handed distinct inputs and want to show the outputs are distinct.

\begin{definition}[surjective]
$f : A \to B$ is \textbf{surjective} if, for every $b \in B$, there exists at least one $a \in A$ with $f(a) = b$. Equivalently: $\mathrm{range}(f) = B$.
\end{definition}

\begin{definition}[bijective]
$f : A \to B$ is \textbf{bijective} if it is both injective and surjective. A bijective function has a \textbf{inverse} $f^{-1} : B \to A$ satisfying $f^{-1}(f(a)) = a$ for every $a$ and $f(f^{-1}(b)) = b$ for every $b$.
\end{definition}

\begin{mistake}
``$f^{-1}(x) = 1/f(x)$.'' No. The superscript $-1$ on a function means the \emph{inverse function}, not the reciprocal. If $f(x) = 2x + 1$, then $f^{-1}(x) = (x - 1)/2$, which is quite different from $1/(2x + 1)$.

Dawkins flags this pointedly on his inverse-functions page: \textit{``Do not confuse the two!''} The notation is confusing enough that mathematicians wrote an explicit warning about it.
\end{mistake}

\subsection*{Arrow diagrams (after MathIsFun)}

Draw the domain $A$ as a blob of dots on the left, the codomain $B$ as a blob of dots on the right, and an arrow from each element of $A$ to the element of $B$ it maps to.

\begin{center}
\begin{tikzpicture}[>=latex,scale=1.0]
  %% injective but not surjective
  \begin{scope}[shift={(0,0)}]
    \node at (-0.5, 2.5) {\textbf{injective, not surjective}};
    \foreach \y in {0,1,2} \fill (-1,\y) circle (2pt);
    \foreach \y in {-0.5,0.5,1.5,2.5} \fill (1,\y) circle (2pt);
    \draw[->] (-1,0) -- (1,-0.5);
    \draw[->] (-1,1) -- (1,0.5);
    \draw[->] (-1,2) -- (1,1.5);
    \node at (-1,-0.7) {$A$};
    \node at (1,-0.9) {$B$};
  \end{scope}
  %% surjective but not injective
  \begin{scope}[shift={(5,0)}]
    \node at (-0.5, 2.5) {\textbf{surjective, not injective}};
    \foreach \y in {0,0.7,1.4,2.1} \fill (-1,\y) circle (2pt);
    \foreach \y in {0.3,1.1,1.8} \fill (1,\y) circle (2pt);
    \draw[->] (-1,0) -- (1,0.3);
    \draw[->] (-1,0.7) -- (1,0.3);
    \draw[->] (-1,1.4) -- (1,1.1);
    \draw[->] (-1,2.1) -- (1,1.8);
    \node at (-1,-0.2) {$A$};
    \node at (1,-0.2) {$B$};
  \end{scope}
  %% bijective
  \begin{scope}[shift={(10,0)}]
    \node at (-0.5, 2.5) {\textbf{bijective}};
    \foreach \y in {0,1,2} \fill (-1,\y) circle (2pt);
    \foreach \y in {0,1,2} \fill (1,\y) circle (2pt);
    \draw[->] (-1,0) -- (1,1);
    \draw[->] (-1,1) -- (1,0);
    \draw[->] (-1,2) -- (1,2);
    \node at (-1,-0.4) {$A$};
    \node at (1,-0.4) {$B$};
  \end{scope}
\end{tikzpicture}
\end{center}

\textbf{Injective:} no two arrows end at the same dot in $B$. (Some dots in $B$ may be unhit.)\\
\textbf{Surjective:} every dot in $B$ has at least one arrow ending at it. (Two arrows may end at the same dot.)\\
\textbf{Bijective:} every dot in $B$ has exactly one arrow ending at it. A perfect pairing.

\subsection*{Worked examples}

\textbf{Example 1.} $f : \mathbb{R} \to \mathbb{R}$ defined by $f(x) = x^2$. \\
Injective? No. $f(2) = 4$ and $f(-2) = 4$. Two distinct inputs, one output. Violates the rule.\\
Surjective? No. The range is $[0, \infty)$. No real input produces $-1$. \\
Bijective? No (failing either injectivity or surjectivity is enough to fail bijectivity).

\textbf{Example 2.} $f : [0, \infty) \to [0, \infty)$ defined by $f(x) = x^2$. Same formula, smaller domain and codomain. \\
Injective? Yes. On $[0, \infty)$, $x^2$ is strictly increasing: if $0 \leq x_1 < x_2$ then $x_1^2 < x_2^2$. Distinct inputs produce distinct outputs. \\
Surjective? Yes. For every $y \geq 0$, the input $x = \sqrt{y}$ is in $[0, \infty)$ and produces $f(\sqrt{y}) = y$. \\
Bijective? Yes. The inverse is $f^{-1}(y) = \sqrt{y}$.

The same formula flipped classification. Restricting the domain and codomain converted a messy non-bijection into a clean bijection. This is why declaring domain and codomain is not cosmetic.

\textbf{Example 3.} $f : \mathbb{R} \to \mathbb{R}$ defined by $f(x) = x^3$. \\
Injective? Yes. If $x_1^3 = x_2^3$, then $x_1 = x_2$ (taking real cube roots). \\
Surjective? Yes. Every real $y$ has a real cube root $\sqrt[3]{y}$, and $f(\sqrt[3]{y}) = y$. \\
Bijective? Yes. The inverse is $f^{-1}(y) = \sqrt[3]{y}$.

Compare Example~1 and Example~3. Structurally they look identical: a single-formula map $\mathbb{R} \to \mathbb{R}$. But $x^2$ fails on both counts and $x^3$ succeeds on both. The reason is that $x^3$ is strictly increasing on all of $\mathbb{R}$, while $x^2$ is decreasing on the negatives and increasing on the positives.

%% ============================================================
%%  WORKED EXAMPLE
%% ============================================================

\section{Worked example: counting arguments as a function}

This example ties together Sections~1 (numbers as counts), 5 (functions), and 6 (domain and codomain). It is worth seeing end to end.

\textbf{Problem.} Given $5$ people, how many ways are there to pick $2$ of them to share a taxi?

\textbf{Step 1. Answer the question concretely.} Label the people $1, 2, 3, 4, 5$ and list every distinct pair:
\[
  \{1,2\},\, \{1,3\},\, \{1,4\},\, \{1,5\},\, \{2,3\},\, \{2,4\},\, \{2,5\},\, \{3,4\},\, \{3,5\},\, \{4,5\}.
\]
Count: $10$. (The order within a pair does not matter; $\{1,2\}$ and $\{2,1\}$ are the same pair.)

\textbf{Step 2. Label the inputs.} The question has two inputs: the total number of people ($n = 5$) and the size of the subset we pick ($k = 2$). The answer depends on both; change $n$ or $k$ and you get a different count.

\textbf{Step 3. Re-state the process as a function.} Define
\[
  C \colon \{(n, k) \in \mathbb{N} \times \mathbb{N} \,:\, 0 \leq k \leq n\} \;\longrightarrow\; \mathbb{N}
\]
by the rule ``$C(n, k)$ is the number of $k$-element subsets of an $n$-element set.'' Then $C(5, 2) = 10$.

Notice what the function's domain and codomain are:
\begin{itemize}
  \item The \textbf{domain} is a subset of $\mathbb{N} \times \mathbb{N}$ (we require $0 \leq k \leq n$; you cannot pick more people than you have, and you cannot pick a negative number of people).
  \item The \textbf{codomain} is $\mathbb{N}$. Counts are natural numbers.
\end{itemize}

\textbf{Step 4. Invoke Section~1.} The output $C(5, 2) = 10$ is a \emph{count}. It is a mathematical number. We can add it (total pairs across two groups), compare it (this configuration has more pairs than that one), and use it arithmetically. This is why it makes sense to talk about the output of $C$ as a number at all. If the output were a label (like a phone number), it would not participate this way.

\textbf{Step 5. (Stated without derivation; full derivation in Chapter~7 or a dedicated combinatorics chapter.)} There is a closed-form formula:
\[
  C(n, k) \;=\; \frac{n!}{k!\,(n-k)!},
\]
where $n! = n \cdot (n-1) \cdot (n-2) \cdots 2 \cdot 1$ is the product of integers from $1$ up to $n$. Check: $C(5, 2) = 5!\,/\,(2!\,3!) = 120/(2 \cdot 6) = 120/12 = 10$. Matches the hand count.

\textbf{Step 6. (Optional bonus question.)} Is $C$ injective as a function $\{(n, k) : 0 \leq k \leq n\} \to \mathbb{N}$? No. $C(2, 1) = 2$ and $C(3, 2) = 3$, so those two are distinct. But also $C(4, 1) = 4$ and $C(4, 3) = 4$: two distinct inputs, same output. So $C$ is not injective. Is it surjective onto $\mathbb{N}$? No (some naturals never appear as binomial coefficients; e.g., $7$ does not appear in the $(n, k)$-table below $n = 7$). This kind of question is exactly what Section~7 prepares you to answer.

\section*{Chapter summary}

If you forget the details, remember these seven points:
\begin{itemize}
  \item counts and measurements are mathematical numbers; labels and positions only borrow digits.
  \item each larger number system is introduced because the smaller one fails to answer some arithmetic question.
  \item the rationals are dense but incomplete; $\sqrt{2}$ exposes one of their holes.
  \item a variable is a symbol ranging over a set, not a box in memory.
  \item a function is a rule assigning exactly one output to each input.
  \item domain, codomain, and range are three different objects and must not be collapsed.
  \item injective, surjective, and bijective are the three core function classifications you will keep using.
\end{itemize}

\newpage

%% ============================================================
%%  QUIZ
%% ============================================================

\section{Quiz}

Ten questions covering all seven sections plus the worked example. Answers in Appendix~A.

\begin{enumerate}
  \item \textbf{(Section~1.)} Which of these are mathematical numbers in the sense of Section~1? Justify each:
    \begin{itemize}[nosep]
      \item ``three apples''
      \item ``phone number 555-1234''
      \item ``the third room on the left''
      \item ``3.5 meters''
    \end{itemize}

  \item \textbf{(Section~2.)} Is subtraction closed over $\mathbb{N}$? If yes, give a brief argument. If no, give a counterexample.

  \item \textbf{(Section~2.)} Give an arithmetic expression that has no answer in $\mathbb{Z}$ but has an answer in $\mathbb{Q}$. What property of $\mathbb{Z}$ does this show is missing?

  \item \textbf{(Section~3.)} Is $\sqrt{2}$ a rational number? Justify using the chapter's proof by contradiction, including the ``$k^2$ even $\Rightarrow k$ even'' lemma.

  \item \textbf{(Section~3.)} Explain briefly in what sense both $\mathbb{N}$ and $\mathbb{R}$ are infinite. Why are their infinities considered different?

  \item \textbf{(Section~4.)} Is $x + 1 = 3x$ an equation or an identity? What about $x + 1 = 1 + x$? Justify each.

  \item \textbf{(Section~4.)} Translate the English sentence ``there is a real number whose square equals $4$'' into symbolic form using $\exists$.

  \item \textbf{(Section~5.)} Is $y^2 = x + 1$ a function? Justify using the chapter's criterion and the vertical line test.

  \item \textbf{(Section~6.)} Find the range of $f : \mathbb{R} \to \mathbb{R}$ defined by $f(x) = 1 / (1 + x^2)$. Is this function surjective?

  \item \textbf{(Section~7.)} Show that $f : \mathbb{R} \to \mathbb{R}$ defined by $f(x) = 2x + 1$ is bijective. Find its inverse.
\end{enumerate}

\newpage

%% ============================================================
%%  APPENDIX A: QUIZ ANSWERS
%% ============================================================

\section*{Appendix A: Quiz answers}

\begin{enumerate}
  \item \textbf{Mathematical numbers.} ``Three apples'' (count, yes). ``Phone number 555-1234'' (label, no). ``The third room on the left'' (position; context-dependent; as a position in an order it is a count of positions, so arguably yes; as a standalone label for a specific room, no). ``3.5 meters'' (measurement, yes).

  \item \textbf{Subtraction closure in $\mathbb{N}$.} Not closed. Counterexample: $3 - 5$ has no answer in $\mathbb{N}$.

  \item \textbf{$\mathbb{Z}$ gap.} $3 \div 7$ has no answer in $\mathbb{Z}$ (no integer multiplied by $7$ gives $3$) but in $\mathbb{Q}$ equals $3/7$. Missing property: $\mathbb{Z}$ is not closed under division by nonzero elements.

  \item \textbf{$\sqrt{2}$ irrationality.} Assume $\sqrt{2} = p/q$ in lowest form with $\gcd(p, q) = 1$. Then $p^2 = 2q^2$, so $p^2$ is even, so (by the lemma) $p$ is even, so $p = 2k$, so $4k^2 = 2q^2$, so $q^2 = 2k^2$, so $q^2$ is even, so $q$ is even. But then $\gcd(p, q) \geq 2$, contradicting the lowest-form assumption. Therefore $\sqrt{2}$ is not rational.

  \item \textbf{Two infinities.} Both $\mathbb{N}$ and $\mathbb{R}$ are infinite. $\mathbb{N}$ is infinite and \emph{listable}: every natural eventually appears in the sequence $0, 1, 2, 3, \ldots$. $\mathbb{R}$ is infinite but \emph{not listable} in that way; Cantor proved that no enumeration covers all the reals. ``Infinite'' is not one size.

  \item \textbf{Equation vs identity.} $x + 1 = 3x$ is an equation: it holds for exactly one value, $x = 1/2$. $x + 1 = 1 + x$ is an identity: it holds for every value of $x$ (it is just the commutativity of addition restated).

  \item \textbf{Quantifier translation.} $\exists x \in \mathbb{R} \colon x^2 = 4$. (Two such $x$ exist, namely $2$ and $-2$; $\exists$ requires at least one, so the assertion is true.)

  \item \textbf{$y^2 = x + 1$ is not a function.} At $x = 3$, $y^2 = 4$ has two solutions $y = +2$ and $y = -2$. The definition requires \emph{exactly one} output per input; two outputs violates the rule. Geometrically, the curve is a sideways parabola; the vertical line $x = 3$ crosses it at two points, $(3, 2)$ and $(3, -2)$, failing the vertical line test.

  \item \textbf{Range of $f(x) = 1/(1 + x^2)$.} $1 + x^2 \geq 1$ for every real $x$, with equality only at $x = 0$. So $0 < f(x) \leq 1$ for every $x$, and $f(0) = 1$ is attained. As $|x| \to \infty$, $f(x) \to 0$ but never reaches it. The range is $(0, 1]$. The codomain $\mathbb{R}$ strictly contains the range, so the function is \textbf{not surjective}. (For example, no input produces $-1$, $0$, or $2$.)

  \item \textbf{$f(x) = 2x + 1$ is bijective.} \\
    \emph{Injective:} suppose $f(a) = f(b)$. Then $2a + 1 = 2b + 1$, so $2a = 2b$, so $a = b$. Distinct inputs produce distinct outputs. \\
    \emph{Surjective:} given any $y \in \mathbb{R}$, the input $x = (y - 1)/2$ satisfies $f(x) = 2 \cdot (y - 1)/2 + 1 = y$. Every element of the codomain is hit. \\
    \emph{Inverse:} $f^{-1}(y) = (y - 1)/2$. Check: $f^{-1}(f(x)) = (2x + 1 - 1)/2 = x$, and $f(f^{-1}(y)) = 2 \cdot (y - 1)/2 + 1 = y$. Both round-trips restore the original.
\end{enumerate}

%% ============================================================
%%  SOURCES
%% ============================================================

\section*{Sources and further reading}

\begin{itemize}
  \item Rosen, K.\ \textit{Discrete Mathematics and Its Applications}, 8th ed.\ McGraw-Hill, 2019.\ Section 2.3 --- the function definition used in Section~5.
  \item Axler, S.\ \textit{Linear Algebra Done Right}, 4th ed.\ Springer, 2024.\ Section 0.A and Chapter~3 --- the three-object naming (domain, codomain, range) used in Section~6.
  \item Stewart, J.\ \textit{Calculus: Early Transcendentals}, 8th ed.\ Cengage.\ Section 1.1 --- the four-way representation taxonomy used in Section~5.
  \item Lang, S.\ \textit{Basic Mathematics}.\ Springer, reprint 1988 (original Addison-Wesley 1971).\ Chapter~1 \S5, Theorem~4 --- the $\sqrt{2}$ irrationality proof adapted in Section~3.
  \item Spivak, M.\ \textit{Calculus}, 4th ed.\ Publish or Perish, 2008.\ Chapters~1--2 --- the P1--P13 axiom framing for $\mathbb{R}$ and the set-$C$ ``hole in $\mathbb{Q}$'' argument used in Section~3.
  \item Dawkins, P.\ \textit{Paul's Online Math Notes}, \texttt{tutorial.math.lamar.edu}.\ Function-definition and inverse-function pages --- the minimal-pair counterexamples and notation warnings used throughout.
  \item MathIsFun, ``Injective, Surjective and Bijective'' --- the arrow-diagram pedagogy and naming-collision warning used in Section~7.
\end{itemize}

\end{document}
